\documentclass[twoside,11pt]{article}

\usepackage{blindtext}
\usepackage{amsmath}

% Any additional packages needed should be included after jmlr2e.
% Note that jmlr2e.sty includes epsfig, amssymb, natbib and graphicx,
% and defines many common macros, such as 'proof' and 'example'.
%
% It also sets the bibliographystyle to plainnat; for more information on
% natbib citation styles, see the natbib documentation, a copy of which
% is archived at http://www.jmlr.org/format/natbib.pdf

% Available options for package jmlr2e are:
%
%   - abbrvbib : use abbrvnat for the bibliography style
%   - nohyperref : do not load the hyperref package
%   - preprint : remove JMLR specific information from the template,
%         useful for example for posting to preprint servers.
%
% Example of using the package with custom options:
%
% \usepackage[abbrvbib, preprint]{jmlr2e}

\usepackage{jmlr2e}
\hypersetup{colorlinks=true, urlcolor=blue, citecolor=blue, linkcolor=blue}

% Definitions of handy macros can go here

\newcommand{\dataset}{{\cal D}}
\newcommand{\fracpartial}[2]{\frac{\partial #1}{\partial  #2}}

% Heading arguments are {volume}{year}{pages}{date submitted}{date published}{paper id}{author-full-names}

\usepackage{lastpage}
\jmlrheading{26}{2026}{1-\pageref{LastPage}}{4/26; Revised 5/26}{9/26}{21-0000}{Aksel Kretsinger-Walters}

% Short headings should be running head and authors last names

\ShortHeadings{A Causal Evaluation of Timeouts in the NBA}{Aksel Kretsinger-Walters}
\firstpageno{1}

\begin{document}

\title{A Causal Evaluation of Timeouts in the NBA}

\author{\name Aksel Kretsinger-Walters \email adk2164@columbia.edu \\
		\addr Department of Computer Science\\
		Columbia University\\
New York, NY 10027, USA}

\editor{David M. Blei}

\maketitle

\begin{abstract}%   <- trailing '%' for backward compatibility of .sty file
		There is a long-standing belief in the NBA that calling a timeout during an opponent's scoring run can
		disrupt their momentum and improve the calling team's chances of winning. However, the empirical evidence
		in support of this belief is limited and inconclusive. In this paper, we use play-by-play data from the
		2021-2026 NBA seasons to conduct a causal analysis of the effect of timeouts on team performance.
\end{abstract}

\begin{keywords}
		Causal inference, sports analytics, NBA, timeouts, momentum, difference-in-differences, synthetic control
\end{keywords}

% The TODO comments are broadly split into latex modifications & experiment/code modifications. We should first
% identify which latex modifications can be done without any code changes (ie are indepdenent of code).
% We should then run through the bottom section of code TODOs and save the updated results
% Then, we update all the remaining latex TODOs that depend on those results
% Finally, we comb through the latex report to find and fix any stale values that weren't explicitly flagged
% with a TODO comment

\section{Significance}
% Why is this problem important? To whom?

The definition and nature of streaks, runs, and momentum itself has been a point of interest and debate for
sports fans and analysts for decades. In addition to the papers cited in our related works, there's a large
body of research on the existence and nature of momentum in sports, both on the individual level (e.g. the
``hot-hand'' effect) and the team level.

On an academic level, the question of whether timeouts disrupt momentum is a interesting case study in causal
inference in a complex, dynamic system. It involves issues of selection bias (coaches call timeouts in
response to game state), confounding (many factors influence scoring dynamics), and the challenge of defining
and measuring momentum itself. The conflicting conclusions in the existing literature suggest that this is a
non-trivial question.

The significance of this question extends beyond academic curiosity. If timeouts do have a causal effect on
disrupting momentum, it would validate a common coaching strategy and could influence how teams manage their
timeouts. There's also opportunity for further discovery: how does the relationship change conditioned on
player substitutions, home-away status, quarter, regular vs post season, team age and
maturity, duration and intensity of the run, etc.? There are many more such ablations to consider. Conversely,
if timeouts are found to have no effect or even a negative effect, it could lead to a reevaluation of how
coaches use their timeouts and potentially change the way teams approach in-game strategy.

\section{Related Works}

The question of whether timeouts actually disrupt momentum has been studied from several angles, with notably
conflicting conclusions. \citet{assis2020stop} apply a formal causal framework to NBA play-by-play data,
using matching techniques informed by a causal graph to estimate the counterfactual effect of a timeout. They
conclude that timeouts have \emph{no} effect on team performance, and argue that the apparent rebound
observed after a timeout is simply regression to the mean: because timeouts tend to be called when the
opposing team is scoring at an unusually high rate, any subsequent slowdown reflects the natural tendency of
scoring to return to its average rather than a causal effect of the stoppage itself.

\citet{gibbs2021causal} revisit the same question but restrict attention to the specific scenario of an
opposing team's scoring run, where the strategic stakes of the timeout decision are highest. Working within
the Rubin causal model, they carefully define units under SUTVA and introduce an interpretable outcome based
on the score difference to capture broad changes in scoring dynamics. In contrast to \citet{assis2020stop},
they find that calling a timeout during an opposing run is \emph{slightly disadvantageous} on average, and
further show that the magnitude of this effect varies meaningfully by franchise. Their result suggests that
comebacks after runs happen frequently enough to create the illusion of timeout efficacy, but that the
intervention itself does not help.

\citet{weimer2023causal} take a different identification strategy by exploiting mandatory television
timeouts, which are triggered at predetermined points in the game and are therefore plausibly exogenous with
respect to in-game momentum. Using TV timeouts as a natural experiment that mimics random assignment of a
pause, they report an \emph{11.2\% decline} in the scoring rate of the team carrying momentum immediately
after a TV timeout, and they show this effect is robust to run size, substitutions, and game context. This
stands in sharp contrast to the null result of \citet{assis2020stop} and the mildly negative result of
\citet{gibbs2021causal}: by removing coach selection bias entirely, \citet{weimer2023causal} argue that
team-level momentum genuinely exists and that pauses in play do interrupt it.

\citet{qiu2025interrupt} approach the problem from a sports-psychology perspective rather than a
causal-inference one, analyzing 4{,}051 timeouts across 1{,}235 elite basketball games using linear mixed
models to test pre- versus post-timeout momentum over short (48s), medium (96s), and long (144s) windows.
They find that timeouts significantly \emph{increase} momentum for teams that are trailing or otherwise
disadvantaged, while tending to \emph{decrease} momentum for teams that are already ahead, with contextual
factors such as game period and opponent strength substantially moderating the effect. Unlike the three
causal papers above, which target an average treatment effect, their results emphasize heterogeneity:
timeouts are neither uniformly helpful nor uniformly useless, and their efficacy is highly conditional on
game state. Together these four studies motivate our own analysis, which aims to reconcile the tension
between the exogenous-timeout evidence for a real momentum effect and the selection-bias-corrected evidence against one.

% Acknowledgements and Disclosure of Funding should go at the end, before appendices and references
\section{Methods}

\subsection{Data}

We build our analysis on the cdnnba enriched play-by-play data for the 2020--2025 NBA seasons (regular
season and playoffs), loaded via our own \texttt{CDNNBAMemoPL.load\_all()} pipeline. After cleaning, the
spine contains 4{,}183{,}347 play-by-play rows spanning 7{,}457 unique games and 1{,}480{,}697 possessions,
of which 83{,}662 are coach-called timeouts and 14{,}561 are inferred mandatory television timeouts.

\subsection{Timeout classification}

We partition stoppage events into three mutually exclusive groups. \textbf{Endogenous} timeouts are
coach-initiated stoppages, identified by \texttt{actionType == timeout} and \texttt{subType} in
\{\texttt{full}, \texttt{challenge}\}. These are endogenous in the causal sense because coaches decide when to
call them in response to game state. \textbf{Exogenous} timeouts are the mandatory television timeouts
specified by the NBA rulebook at roughly the 6:59 and 2:59 marks of each period. The cdnnba feed does not
directly label TV timeouts, so we inject them using a rulebook-based inference procedure and assign them
\texttt{subType == official\_inferred}. (TODO update the efficacy of this procedure) We validated this
procedure against the explicitly labeled 2013--2016
nbastatsv3 feed and obtained an F1 of $0.88$--$0.92$, indicating good but imperfect recovery of mandatory
timeouts. \textbf{Control} events are dead-ball moments with no timeout occurring during an active scoring
run; to avoid over-representing long runs, we subsample at most one control per run segment.

The endogenous/exogenous distinction is causal rather than mechanical. NBA rules allow substitutions and
strategy huddles during \emph{any} timeout (full, short, or mandatory TV), so the opportunity to alter play
after the break is identical in both cases. What separates the two types is how the trigger is selected:
coaches call endogenous timeouts in response to game state (introducing selection bias on adverse momentum),
whereas mandatory TV timeouts fire at fixed clock positions and are therefore plausibly exogenous
conditional on the clock.

\subsection{Run and recovery definitions}

A \emph{scoring run} at play $t$ is defined as $|\texttt{streak}(t)| \geq s$, where $\texttt{streak}(t)$ is
the cumulative net score swing within the current scoring segment, computed by resetting the counter whenever
the scoring team flips. We refer to the team being outscored as the \emph{suffering team} and the team
scoring as the \emph{running team} or alternatively the \emph{streaking team}.

Our primary outcome is the \emph{recovery} metric: the signed lead change from the suffering team's
perspective over a forward window of $n$ minutes. If $\texttt{lead}(t)$ is the home-minus-away margin at play
$t$, we define
\[
		\texttt{recovery}(t, n) = -\operatorname{sign}(\texttt{streak}(t)) \cdot \bigl(\texttt{lead}(t + n\text{
		min}) - \texttt{lead}(t)\bigr),
\]
so that positive values mean the suffering team clawed points back during the window. As a secondary outcome
we also report $\texttt{running\_team\_pts}$, the raw points scored by the running team over the forward
window, which matches the dependent variable used by \citet{weimer2023causal}.

\subsection{Between-group estimators}

As a first pass we compute group means of \texttt{recovery} conditional on run size $s$ and forward window
$n$ (measured in minutes), and report differences versus the control group using Welch's $t$-tests. We sweep
$s \in \{4, 5, 6, 7,
8, 10, 12, 15\}$ and $n \in \{0.5, 1, 1.5, 2, 3, 4, 5\}$ minutes, and also stratify by quarter, regular season
vs.\ playoffs, and absolute score margin.

\subsection{Within-game matched-twin causal design}\label{sec:matched-twin}

\paragraph{Motivation.}
The between-group estimator is vulnerable to game-level confounding: coaches call timeouts in systematically
different games than the ones where control events are drawn, and mandatory TV timeouts cluster at
predictable game-clock positions (around the 6:59 and 2:59 marks of each period) rather than uniformly
across all game contexts. Comparing the timeout pool against a global control mean therefore absorbs game-
and clock-level selection into the baseline. To remove this bias we use a \emph{within-game matched-twin}
design that holds team composition, pace, coaching, and game-level momentum approximately fixed by
construction.

\paragraph{Method.}
We first build a candidate pool by classifying every spine event during a run of $|\texttt{streak}| \geq s$
into one of three groups: \textbf{treated-endogenous} (a coach \texttt{full} or \texttt{challenge} timeout
called by the suffering team), \textbf{treated-exogenous} (an \texttt{official\_inferred} TV timeout or a
generic dead-ball \texttt{stoppage}), or \textbf{control} (any spine event during the same run that is not
		classified as treated, with at most one control sampled per contiguous run segment to avoid over-counting
long runs). For each treated event we then search the control pool \emph{within the same game} for a
candidate that satisfies all of: the same period; the same sign of $\texttt{streak}$ (so the run is against
the same team); $|\Delta\,\texttt{streak}| \leq 2$ (comparable run magnitude);
$|\Delta\,\texttt{seconds\_remaining}| \leq 120$\,s (comparable clock position within the period); and
$|\Delta\,\texttt{suffering\_margin}| \leq 3$, where the suffering margin
\[
		\texttt{suffering\_margin}(t) \;=\; -\operatorname{sign}(\texttt{streak}(t)) \cdot \texttt{lead}(t)
\]
is the score margin from the suffering team's perspective (positive when the suffering team is still ahead
despite the run against it). When multiple candidate controls satisfy the constraints, we select the one
minimising the weighted Manhattan distance
\[
		d(\text{treated},\, \text{control}) \;=\; 2\,|\Delta\,\texttt{streak}|
		\;+\; \frac{|\Delta\,\texttt{seconds\_remaining}|}{30}
		\;+\; 3\,|\Delta\,\texttt{suffering\_margin}|,
\]
which trades a ${\sim}30$-second clock gap for one point of streak magnitude or one point of score margin.
Each treated event yields at most one pair; the same control row may serve as the matched twin for multiple
treated events. The per-pair causal estimate is the difference in $\texttt{recovery}$ between the treated
event and its matched twin, and we aggregate across pairs with a paired $t$-test.

\subsection{Moderator analyses}

To probe heterogeneity we repeat the between-group and matched-twin analyses within slices defined by
(i) timeout subtype (\texttt{coach\_full}, \texttt{coach\_challenge}, \texttt{official\_inferred},
\texttt{stoppage}); (ii) game phase and clutch status (last five minutes of Q4 with margin $\leq 5$);
(iii) season-long team quality gap, using each team's \texttt{NET\_RATING} from \texttt{player\_advanced\_stats};
(iv) the number of substitutions within a $\pm 30$s game-clock window around the event, bucketed as
$\{0,\, 1\text{--}2,\, 3{+}\}$; and (v) the running team's in-game cumulative points-per-possession (PPP) at
the time of the event, as a proxy for short-term shooting rhythm.

\section{Evaluation}

\subsection{Between-group results}

The simple between-group comparison essentially reproduces the weak-to-null effects reported in the prior
literature. Across run sizes $s \in \{5,\dots,10\}$ and forward windows $n \in \{1,\dots,5\}$ minutes, the
endogenous-vs-control gap in \texttt{recovery} is never statistically significant (all $|\Delta| < 0.1$ pts,
$p > 0.05$ in 15 of 16 cells). Exogenous timeouts, by contrast, look significantly \emph{negative}, with
$\Delta \approx -0.25$ to $-0.45$ pts at the 3-minute window depending on run size ($p < 0.001$). Sweeping
the forward window at fixed run size $s = 6$ shows the exogenous penalty growing from $-0.20$ at 30 seconds
to a plateau around $-0.27$ at 3 minutes, while the endogenous gap hovers near zero throughout. At face value
this matches the ``timeouts don't help (and mandatory TV breaks hurt)'' consensus. As a secondary check we
also computed $\texttt{running\_team\_pts}$ in the style of \citet{weimer2023causal} and found a
\emph{positive} exogenous-minus-control gap of $+2.8\%$ to $+14.6\%$ depending on window  --  the opposite sign
from Weimer's reported $-11.2\%$. This is internally consistent with our suffering-team analysis (if the
suffering team recovers less, the running team must score relatively more) and likely reflects the post-2017
mandatory-timeout rule changes together with the inferred-rather-than-labeled nature of our TV timeouts.

To avoid confusion we'll restate our conclusion for this result: coach timeouts appear to perform worse
than random stoppages of play and exogenous (TV) timeouts. We believe that this is \textbf{not} a causal
result, but instead a consequence of the adverse situations during which coaches predominantly call
timeouts. TODO: update this conclusion once the matched-twin results are in, and make sure to clarify the
distinction between the between-group and within-game results in the final report.

\subsection{Within-game matched-twin results}
TODO: update the numbers and significance levels below once the new matched-twin experiments are done, and
update the interpretation accordingly. We should also add a sentence or two about the challenge results,
which are currently missing from the report.

The matched-twin design of \S\ref{sec:matched-twin} -- our cleanest causal test -- reverses the headline
conclusion of the between-group analysis. With $s = 6$ and a 3-minute window we obtain 8{,}906 matched
endogenous pairs and 2{,}667 matched exogenous pairs (each treated event matched to its closest in-game
		control by the weighted distance $d$; control rows are reused across treated rows when they are the best
match for multiple). The paired mean differences are \textbf{$+0.397$ pts for endogenous timeouts}
($t = 27.9$, $p < 0.0001$) and \textbf{$+0.455$ pts for exogenous timeouts} ($t = 11.8$, $p < 0.0001$).
Both effects are highly significant and both favor the suffering team; the exogenous effect is in fact
slightly \emph{larger} than the endogenous effect, which is consistent with coaches selectively burning
timeouts in especially adverse game states where recovery is hard even with help (see
\hyperref[sec:future]{Future Work} on improving the pairing). Breaking the treated population into fine
subtypes, every subtype shows the same sign and roughly the same magnitude: \texttt{coach\_full} $+0.396$
***, \texttt{coach\_challenge} $+0.714$ * (small sample, $n = 35$), \texttt{official\_inferred} $+0.449$
***, and generic \texttt{stoppage} events $+0.458$ ***. The reversal relative to the between-group results
is driven by game-level selection: coach timeouts concentrate in games where the suffering team historically
claws back less of its deficit during opponent runs (low-recovery games), so comparing the endogenous mean
to a global control mean absorbs that selection into the baseline and erases the signal. Matching within
game cancels this bias by construction. TODO: what does this mean? A coarser within-game \emph{mean}
comparison -- in which each
treated event is compared against the average $\texttt{recovery}$ of all controls in the same game, rather
than against a single distance-minimising twin -- gives a consistent $+0.35$ pts effect for endogenous
timeouts but shrinks the exogenous effect to near zero. The drop suggests that tight clock-position
matching is more load-bearing for exogenous events (whose triggers are themselves clock-locked) than for
endogenous ones.

We also offer a hypothesis for the exaggerated impact that coaching challenges demostrate in the data. A
successful coaching challenge offers the calling team a break in play, a psychological boost, and tangible
in-game benefits (e.g. Free Throws, Possession, etc). If we assume that coaches are skilled at challenging
calls, it would follow to reason that this category of game stoppage would show larger effects. The
assumption on successful coaching challenges is intuitive: coaches are backed by large analytics and rules
experts teams, and the finite number of challenges (1 per game, with a second awarded if the first challenge
is successful) means that coaches must use them judiciously. (TODO: re-run notebooks and figure out how the
fail / success split looks )

\subsection{Mechanism and moderators}

Once we have a signal, we can ask where it comes from. All moderator analyses in this section use a
between-group Welch test of treated vs control within each stratum -- the same estimator as \S4.1, not the
paired matched-twin of \S4.2. The magnitudes below are therefore directly comparable to the near-zero
between-group baseline and reflect the residual signal that survives the selection bias after stratifying,
not the underlying causal effect size; we use the coarser estimator here to keep per-stratum sample sizes
large enough for a clean test.

Stratifying by the number of substitutions in a $\pm 30$s window around the event reveals that the
endogenous effect is concentrated entirely in timeouts accompanied by personnel changes: with 0
substitutions the endogenous-control gap is $-0.05$ (n.s.), with 1--2 substitutions it is $+0.10$ (n.s.),
and with 3 or more substitutions it rises to $+0.21$ ($p < 0.001$). The exogenous penalty is
correspondingly largest ($-0.39$, $p < 0.001$) when no subs occur and vanishes when many do. This directly
confirms the mechanism proposed by \citet{weimer2023causal}: the value of a timeout is mediated by
substitutions rather than by the pause itself. Conditioning on team quality
via season-long \texttt{NET\_RATING}, we find that the endogenous benefit is strongest in evenly matched
games ($+0.23$ pts, $p < 0.05$) and washes out in lopsided matchups where the talent gap dominates the
trajectory. Finally, stratifying by the running team's in-game PPP we find that when the opponent is shooting
unusually efficiently (PPP $\geq 1.30$), both effects are amplified and point in opposite directions:
endogenous timeouts produce $+0.24$ pts of extra recovery ($p < 0.05$) while exogenous timeouts produce
$-0.62$ pts ($p < 0.001$). This is the clearest evidence in our data of a genuine momentum effect  --  a coach
timeout with substitutions disrupts a hot-shooting opponent, while a pure TV break effectively gives that
opponent a rest and lets the run continue.

\subsection{Summary}

Our headline finding is that timeouts do causally improve the suffering team's short-term recovery  --  by
roughly $+0.4$ points over the following three minutes  --  and that this effect is masked in between-group
analyses by a subtle but severe game-level selection bias. The published consensus of ``no effect'' or
``slightly negative effect'' \citep{assis2020stop,gibbs2021causal} reproduces in our data only when we use a
between-group estimator; the within-game matched-twin estimator, which holds game identity approximately
fixed by construction, produces a strongly positive and highly significant effect for both endogenous and
exogenous stoppages. The benefit is overwhelmingly mediated by substitutions (in agreement with
\citet{weimer2023causal}) and is largest in evenly matched games and against hot-shooting opponents, which
matches both coaching intuition and the conditional-momentum findings of \citet{qiu2025interrupt}. TODO:
update summary once new experiments are done

\section*{Future Work}\label{sec:future}

Several extensions follow naturally from the present design:
\begin{itemize}
		\item \textbf{Better pairing mechanism.} How should we score the quality of a matched twin? The
				current Manhattan-weighted distance $d$ is a reasonable first cut, but there is a sweet-spot
				to characterise between pair similarity (which tightens the estimate) and matched-pair
				sample size (which determines its precision).
		\item \textbf{Alternative outcome metrics.} Beyond $\texttt{recovery}$, we could measure the
				streaking team's points-per-possession or points-per-minute over the forward window, or roll
				either up into a win-probability impact via a calibrated WP model.
		\item \textbf{Causal graphical model.} Distil the moderator findings into an explicit CGM
				(substitutions $\to$ recovery; running-team PPP $\to$ recovery; \texttt{NET\_RATING} gap
				$\to$ recovery) and use it to formalise the identification arguments and check for missing
				confounders.
		\item \textbf{Streaking-team timeouts.} Re-run the analysis with the \emph{streaking} team as the
				calling team. Do they substitute players at a higher rate than baseline, and what mechanism
				would explain a positive or negative effect on the run they are already riding?
		\item \textbf{Additional break types.} Official reviews, injury stoppages, and ejections all
				introduce small exogenous perturbations to play. Do they show the same recovery pattern as
				TV timeouts? Cross-checking against more natural experiments would strengthen the causal
				claim.
\end{itemize}

\section*{Further experimentation and code changes (TODO delete section once complete)}
\begin{itemize}
		\item ``how does the timeout impact for coaching challenges look if we split by successful / failed
				challenges? add the code to the .py
				file, and invoke from the notebook (place the call in the relevant section)''
		\item ``there's also a chance that successful challenges aren't marked as timeouts. take another look at
				the event types and check my
				idea. I think it would also be good to look at the difference in real-world-time (ie 9:43:38.2 pm ->
				9:47:11.9 PM) to indicate
				stoppages in play (in case that reveals where the successful challenges might be hiding). Also I thought
				that there were 35
				matched outcomes (according to the paper) - why do you only se 19 now??''
		\item use real-world-elapsed-time as a feature in TV timeout injection model (retrain and evaluate on nbastatsv3)
		\item reconcile values between section 4.2 and 4.3
		\item ensure that we're only including timeouts called by the suffering team in the endogenous filter
		\item analyze real-world-elapsed-time as a mechanism
		\item team age as a moderator (use average age of players on the court at the time of the timeout). Also
				use the delta between avg age of players on the court before and after the break
		\item Should we be more selecive about including datapoints at the end of the quarters? i.e. never include
				datapoints in the last k (such as 2) minutes of the quarter, and don't include datapoints if t + n
				exceeds the length of the quarter? ie with 3.5 minutes left, include the 3 minute recovery value, but
				exclude the 4 minute window from analysis
		\item update NBA data? (I think this would be expensive and tricky becuase Ted's data is impossilbe to download)
		\item Can we investigate simple substitutions w/out official timeouts? Such as free throws. When else are
				teams allowed to substitute players?
		\item Re-run the study, but replace minutes with \# possessions (make sure this is an even number, to keep
						possessions per team the same. an edge case is a technical foul, which awards the transgressed team a
				free possession. this is infrequent enough that it might not matter)
		\item reconcile with Weimer finding: he states that exogenous breaks cause 11.2 \% drop in scoring by the
				streaking team has. Do we match that?
		\item reconcile challenges in matched experiment (section 4.2) with E15.
		\item re-run experiment post-fixes, flag differences, update conclusion if necessary
\end{itemize}

\acks{N/A}

\section*{LLM Usage}

\texttt{claude\/opus-4.6} was utilized extensively during the implementation of this research project, and we
wanted to acknowledge its contributions.

The author(s) were responsible for project idea generation, sourcing relevant papers and publications,
research roadmap, code scaffolding, and report skeleton.

\texttt{Opus} assisted in the implementation: data preprocessing and cleaning, feature generation,
significance testing, ablation studies, and report completion. All \texttt{Opus} work was closely reviewed
and cross-examined.

% Manual newpage inserted to improve layout of sample file - not
% needed in general before appendices/bibliography.

\newpage

\bibliography{references}

\section*{Appendix}

TODO: add full research-summary.md ablation tables here.

\end{document}
